\chapter{Literature Review and Related Work}
\label{chap:literature_review}

\section{Introduction to the Literature Landscape}
Prognostics and Health Management (PHM) for mission-critical aeromechanical assets has evolved dramatically over the past two decades. The progression from classical physics-based fatigue models to modern data-driven deep learning architectures reflects increasing sensor telemetry availability and algorithmic sophistication. However, the operational adoption of data-driven Remaining Useful Life (RUL) predictors in aviation maintenance scheduling remains severely constrained by two fundamental deficiencies: the absence of rigorous, distribution-free uncertainty quantification, and the disconnect between cycle-level regression accuracy and downstream maintenance decision utility.

This chapter presents a structured synthesis of the literature across five thematic pillars:
\begin{enumerate}
    \item Turbofan Remaining Useful Life (RUL) prediction paradigms;
    \item Uncertainty quantification and conformal calibration;
    \item Operating regime heterogeneity and subgroup reliability;
    \item Maintenance decision theory and conservatism controls;
    \item The specific research gap addressed by the IntelliTwin framework.
\end{enumerate}

\section{Pillar 1: Turbofan Remaining Useful Life Prediction}
The foundational benchmark for data-driven turbofan degradation research is the NASA Commercial Modular Aero-Propulsion System Simulation (C-MAPSS) dataset introduced by Saxena et al.~\cite{saxena2008damage}. C-MAPSS models thermodynamic degradation in commercial turbofans under varying flight profiles, providing run-to-failure trajectories for high-pressure compressor and fan deterioration.

Early algorithmic approaches focused on classical machine learning, such as Support Vector Regression (SVR), Random Forests, and Gradient Boosted Decision Trees (GBDT) \cite{huang2021survey}. With the deep learning paradigm shift, recurrent neural architectures became the prevailing standard. Recurrent architectures, such as Long Short-Term Memory (LSTM) networks and attention mechanisms~\cite{diao2026turbofan}, have been widely applied to capture temporal degradation dynamics in multivariate telemetry alongside convolutional networks~\cite{chen2022data}.

However, across the turbofan prognostics literature, model evaluation has remained almost universally confined to symmetric error metrics, specifically Root Mean Squared Error (RMSE) and Mean Absolute Error (MAE), or ad-hoc asymmetric exponential scoring functions \cite{saxena2008damage}. These point-prediction paradigms fail to provide decision-makers with confidence bounds necessary for safety-critical maintenance dispatch.

\section{Pillar 2: Uncertainty Quantification and Conformal Prediction}
To address the limitations of point estimates, various uncertainty quantification (UQ) methodologies have been explored in PHM:
\begin{itemize}
    \item \textbf{Bayesian and Probabilistic Frameworks}: Predictive uncertainty methods differ in their assumptions and calibration behavior. Their maintenance value must be evaluated through interval reliability and downstream decision outcomes under a common protocol.
    \item \textbf{Quantile Regression}: Estimates conditional quantiles directly via the pinball loss. However, quantile regression provides no finite-sample coverage guarantees and is vulnerable to quantile crossing under non-stationary degradation.
    \item \textbf{Conformal Prediction}: Conformal prediction theory provides a distribution-free uncertainty framework \cite{javanmardi2023conformal}. By evaluating nonconformity scores on a reserved calibration dataset, inductive split-conformal inference constructs prediction intervals targeting valid finite-sample marginal coverage:
    \begin{equation}
    P(Y_{n+1} \in \mathcal{C}(X_{n+1})) \ge 1 - \alpha.
    \end{equation}
\end{itemize}

Recently, Javanmardi et al.~\cite{javanmardi2023conformal} and Yang et al.~\cite{yang2026empirical} explored conformal prediction for industrial prognostics. However, standard split-conformal methods assume exchangeability, which is challenged by non-stationary degradation trajectories and multi-regime flight envelopes.

\section{Pillar 3: Operating Regime Heterogeneity and Subgroup Reliability}
In real-world aviation operations, engines traverse diverse flight regimes (climb, cruise, descent, sea-level take-off). In C-MAPSS sub-datasets FD002 and FD004, six discrete operating regimes induce severe sensor non-stationarity.

When conformal calibration is applied globally across a pooled calibration split, the resulting marginal guarantee $1-\alpha$ guarantees only aggregate coverage across the entire dataset. Lee, Barber, and Willett~\cite{lee2023hierarchical} established hierarchical conformal inference for clustered data, placing $\frac{1}{m+1}$ probability mass at $+\infty$ to target valid coverage for unseen clusters under exchangeable group sampling.

IntelliTwin uses an engine-weighted conformal construction motivated by hierarchical conformal prediction. Formal coverage results require the hierarchical sampling and exchangeability assumptions stated in the underlying theory. Equal engine weighting does not establish these assumptions for ordered, lifetime-stratified degradation trajectories. This study therefore evaluates marginal and subgroup coverage empirically. It does not claim guaranteed coverage for each engine, lifecycle stage, operating regime, complete trajectory, or adaptively selected maintenance time. Nominal 90% coverage is a target, not a guaranteed realized coverage proportion. Maintenance outcomes are evaluated separately under the frozen simulator.

\section{Pillar 4: Maintenance Decisions and Conservatism Controls}
The ultimate objective of prognostic modeling is to inform maintenance scheduling. In operational fleet logistics, an asymmetric cost structure prevails:
\begin{equation}
C_{\text{unplanned}} \gg C_{\text{preventive}}.
\end{equation}
Overestimating RUL allows an engine to operate until catastrophic functional failure, incurring heavy safety penalties, whereas underestimating RUL merely causes premature scheduled replacement, forfeiting a fraction of usable life \cite{saxena2008damage, robinson2026riskaware}.

Cortes-Gomez et al.~\cite{cortesgomez2025utility} introduced utility-directed conformal calibration, while Robinson et al.~\cite{robinson2026riskaware} explored risk-aware RUL bounds. However, an unresolved question in the literature is whether observed maintenance cost reductions stem from genuine uncertainty calibration or merely from trivial over-conservatism (i.e., triggering maintenance earlier regardless of prediction accuracy). Resolving this attribution requires evaluating calibrated intervals against rigorously matched-width heuristic baselines, widened intervals, and fixed safety margins.

\section{Pillar 5: The IntelliTwin Synthesis and Research Gap}
Despite the individual advancements cataloged above, no prior study has simultaneously established:
\begin{enumerate}
    \item Complete engine-level data leakage isolation without cycle-level shuffling;
    \item Condition-stratified conformal calibration with deterministic fallback for rare flight regimes;
    \item Asymmetric nonconformity scaling aligned with maintenance cost structures;
    \item End-to-end discrete-event maintenance simulation evaluating sequential fleet utility;
    \item Attribution testing against matched-conservatism controls under prespecified cluster-bootstrap statistical protocols.
\end{enumerate}
IntelliTwin directly resolves this research gap through an open, reproducible, and prespecified experimental framework version-controlled prior to held-out evaluation.
